\section{Definitions}\label{sec:txid-definitions}

A tracing scheme involves four types of participants. The ledger operator
maintains a publicly accessible, append-only ledger that stores submitted
transactions, and runs the transaction validators that decide, on input of the
ledger, which transactions are valid; we assume the ledger operator is honest,
so that the ledger itself is available and correctly maintained. The
registration authority is honest-but-curious during the protocol: it
authenticates users and issues the credentials required to submit transactions
and enable selective traceability. We additionally assume it does not bias the
public parameters, which in the instantiation of~\Cref{sec:txid-instantiation}
is discharged by deriving all group generators from a public seed rather than
sampling them. Users are potentially malicious, and generate and submit
confidential transactions while retaining fine-grained control over which
auditors are authorised to trace them, and during which epochs. Auditors are
potentially malicious, and may be granted permission to trace the transactions
of specific users during specific epochs.

We consider static corruption throughout: the simulator corrupts users and
auditors of its choice before any call to the ideal functionality is made.

\begin{description}

  \item[$(\crs,\state) \sample \IDgen(\secparam)$ :] The setup algorithm, on
input security parameter $\secparam$, generates the public ledger parameters
$\crs$ together with the initial ledger state $\state$.

  \item[$(\sk,\utk,\regcred) \sample \IDregister(\crs,\state)$ :] The
registration algorithm, run interactively between a user and the registration
authority on input public parameters $\crs$ and state $\state$, outputs user
signing key $\sk$, tracing key $\utk$, and registration credential $\regcred$ to
the user. The instantiation takes $\utk$ to be a tuple.

  \item[$(\PTK,\ek) \sample \IDgrant(\crs,\sk,\utk,\regcred,\epoch)$ :] The
authorisation algorithm, run interactively between a user and an auditor on
input public parameters $\crs$, signing key $\sk$, tracing key $\utk$,
registration credential $\regcred$, and epoch $\epoch$, outputs the user's certified public tracing key
$\PTK$ and the epoch tracing key $\ek$ to the auditor. Only $\ek$ depends on
$\epoch$, so an auditor authorised for several epochs of the same user receives
$\PTK$ once.

  \item[$\tx \sample \IDsubmit(\crs,\sk,\utk,\regcred,\epoch,\payload)$ :] The
submission algorithm, on input public parameters $\crs$, signing key $\sk$,
tracing key $\utk$, registration credential $\regcred$, epoch $\epoch$, and
transaction payload $\payload$, outputs transaction $\tx$, which the user sends to the ledger
operator.

  \item[$\{0,1\} \leftarrow \algo{Verify}(\crs,\state,\tx)$ :] The validation
algorithm, run by the transaction validators on input public parameters $\crs$,
current ledger state $\state$, and transaction $\tx$, returns $1$ if $\tx$ is to
be appended to the ledger and $0$ otherwise. It is non-interactive: a validator
contacts neither the sender, nor the registration authority, nor any auditor.

  \item[$\{\payload\} \leftarrow \IDidentify(\crs,\PTK,\ek,\epoch)$ :] The
tracing algorithm, on input public parameters $\crs$, public tracing key $\PTK$,
epoch tracing key $\ek$, and epoch $\epoch$, returns the set of payloads of the
transactions submitted under $\ek$ during $\epoch$.

\end{description}

\subsection{Ideal functionality $\idfu{\algo{trace}}$}

We define security of tracing through an ideal functionality
$\idfu{\algo{trace}}$ that represents a trusted third party which executes the
operations of the ledger on behalf of system participants, and which by
construction satisfies the three properties set out below.

The ideal functionality captures the properties of tracing: anonymity,
non-repudiation, and permission expiration. Anonymity is enforced by ensuring
that only an auditor holding a valid tracing permission for a specific user and
epoch can de-anonymise that user's transactions during that epoch.
Non-repudiation is achieved by ensuring that every transaction a user submitted
during an epoch can be traced back to them by any auditor holding a valid
permission for that user and epoch. Permission expiration is enforced by
ensuring that tracing permissions are automatically revoked at the end of the
epoch for which they were granted, so an auditor retains no tracing capability
beyond it.

Two things this functionality does \emph{not} capture are worth stating plainly.
First, it models unlinkability of a transaction to its originator, not
confidentiality of the transaction payload: the ledger is public, so
$\idfu{\algo{trace}}$ hands every submitted payload to $\simu$ on every
$\IDsubmit$ call, and what it withholds is the identity of the submitting user.
Confidentiality of the payload is inherited from the underlying transaction
system and is orthogonal to what we prove here. Second, $\IDgrant$ places no
constraint on $\epoch$ relative to $\epoch_{\mathsf{cur}}$. A user may grant a
permission for a future epoch, which is the intended usage, and may also grant
one for an epoch that has already elapsed, which amounts to voluntary
retrospective self-disclosure. What a user cannot do is refuse to be traced in
an epoch for which they have already granted a permission.

\paragraph{Security model.} We work in the standalone model with sequential
composition rather than the universal composability framework. Let
$\textsc{Exec}_{\Pi,\adv}(\secparam,\mathcal{Z})$ denote the joint output of the
honest parties and $\adv$ in a real execution of $\Pi$ with inputs supplied by
environment $\mathcal{Z}$, and let
$\textsc{Exec}_{\idfu{\algo{trace}},\simu}(\secparam,\mathcal{Z})$ denote the
joint output of the honest parties and $\simu$ interacting with
$\idfu{\algo{trace}}$ in the ideal world with the same environment. We say that
$\Pi$ securely realises $\idfu{\algo{trace}}$ if for every PPT adversary $\adv$
there exists a PPT simulator $\simu$ such that for every PPT environment
$\mathcal{Z}$ the two distributions are computationally indistinguishable. The
distinction from universal composability matters because our simulator extracts
witnesses by rewinding the adversary,
via~\Cref{prot:txid-tag-sigma,prot:txid-tx-nizk} and the forking lemma, which is
sound in the standalone setting but not in the UC setting. Straight-line
extraction, and with it UC security, can be recovered by applying the Fischlin
transform to both sigma protocols, or by having the prover verifiably encrypt
its witness under a key in the CRS, at a cost in proof size that we do not pay
here.

$\idfu{\algo{trace}}$ receives inputs from system participants and $\simu$,
computes the corresponding outputs, and delivers them to the authorised parties.
It exposes interfaces through which $\simu$ and honest participants submit
inputs, and additionally maintains a current epoch counter that advances by one
whenever the system moves to the next epoch. The formal description is given
in~\Cref{fig:ftxid}. We summarise each interface below.

\paragraph{$\IDregister$.} A user calls $\IDregister$ to join the system. If the
registration authority approves and the user has not previously registered,
registration succeeds.

\paragraph{$\IDsubmit$.} A registered user calls $\IDsubmit$ to submit up to
$\Max$ transactions during the current epoch. Each successful call records the
payload together with the identity of the user and the current epoch. $\simu$
learns the payload on every call, since the ledger is public, along with the
user's identity if the user is corrupt. The functionality subsumes
$\algo{Verify}$: a call that passes its checks corresponds to a transaction that
the validators accept.

\paragraph{$\IDgrant$.} A user calls $\IDgrant$ to authorise an auditor to trace
their transactions for a specific epoch. $\idfu{\algo{trace}}$ records that the
auditor is authorised to trace the user's transactions for that epoch.  Once the
epoch elapses, the authorisation lapses with it.

\paragraph{$\IDidentify$.} An auditor who has been granted permission for a
given user and epoch calls $\IDidentify$ to obtain the payloads of that user's
transactions during the epoch. $\simu$ additionally learns the traced payloads
if the auditor is corrupt.

\begin{figure} \centering \hspace*{-0.05\textwidth}
\begin{minipage}{1.1\textwidth}

    \begin{framed} \scriptsize

      \setlist[enumerate]{ topsep=2pt, partopsep=0pt, parsep=0pt, itemsep=1pt }

      $\idfu{\algo{trace}}$ maintains tables $\reg$, $\txns$, and $\perm$ to
record user activity: $\reg$ tracks registration status, $\txns$ tracks the
payloads users have submitted, and $\perm$ tracks the tracing permissions
granted to auditors. It also maintains a current epoch counter
$\epoch_{\mathsf{cur}}$, initialised to $1$. We consider a static corruption
model, where $\simu$ corrupts users and auditors of its choice before any call
to $\idfu{\algo{trace}}$ is made. $\Cor$ is the set of corrupt participants.

      \begin{multicols}{2}

	{$\IDregister$.} On input register from the user $\user$, send
$(\IDregister, \user)$ to the registration authority and wait for confirmation.
On receiving confirmation, do the following:

	\begin{enumerate}

	  \item If $\reg[\user] \neq \bot$, then stop. (A user can register at
most once.)

	  \item Send $(\IDregister, \user)$ to $\simu$ and wait for the
go-ahead. On receiving it, continue.

	  \item Set $\reg[\user] \leftarrow \mathsf{true}$.

	  \item Return $\IDregisterend$ to $\user$ and $\simu$.

	\end{enumerate}

	{$\IDsubmit$.} On input $(\IDsubmit, \payload, \epoch)$ from $\user$,
do:

	\begin{enumerate}

	  \item If $\reg[\user] = \bot$, then stop. (Only registered users can
submit transactions that will be accepted.)

	  \item If $\epoch \neq \epoch_{\mathsf{cur}}$, then stop.

	  \item If $|\txns[(\user,\epoch_{\mathsf{cur}})]| + 1 > \Max$, then
stop. (Users can submit at most $\Max$ transactions per epoch.)

	  \item If $\user \in \Cor$, then $\mathsf{origin} \leftarrow \user$,
else $\mathsf{origin} \leftarrow \varnothing$.

	  \item Send $(\IDsubmit, \payload, \epoch, \mathsf{origin})$ to
$\simu$, and wait for the go-ahead. (The ledger is public, so $\simu$ always
learns the payload; if the user is corrupt, $\simu$ additionally learns the
origin of the transaction.)

	  \item Set $\txns[(\user,\epoch_{\mathsf{cur}})] \leftarrow
\txns[(\user,\epoch_{\mathsf{cur}})] \cup \{\payload\}$.

	  \item Return $\IDsubmitend$ to $\user$ and $\simu$.

	\end{enumerate}

	{$\IDgrant$.} On input $(\IDgrant, \auditor, \epoch)$ from $\user$, send
$(\IDgrant, \user, \epoch)$ to auditor $\auditor$ and wait for confirmation. On
receiving confirmation, do:

	\begin{enumerate}

	  \item If $\reg[\user] = \bot$, then stop.

	  \item Send $(\IDgrant, \auditor, \user, \epoch)$ to $\simu$ and wait
for the go-ahead. On receiving it, do:

	    \begin{enumerate}

	      \item Set $\perm[(\auditor,\user,\epoch)] \leftarrow
\mathsf{true}$.

	      \item Return $\IDgrantend$ to $\user$, $\auditor$, and $\simu$.

	    \end{enumerate}

	\end{enumerate}

	{$\IDidentify$.} On input $(\IDidentify, \user, \epoch)$ from auditor
$\auditor$, do:

	\begin{enumerate}

	  \item If $\perm[(\auditor,\user,\epoch)] \neq \mathsf{true}$, then
stop. (Only authorised auditors can trace the transactions of a user.)

	  \item Otherwise, set $T \leftarrow \txns[(\user,\epoch)]$.
($\auditor$ traces every transaction $\user$ submitted during $\epoch$:
non-repudiation.)

	  \item If $\auditor \in \Cor$, then $T^\star \leftarrow T$, else
$T^\star \leftarrow \varnothing$.

	  \item Return $(\IDidentifyend, T)$ to $\auditor$ and $(\IDidentifyend,
T^\star)$ to $\simu$. (If the auditor is corrupt, $\simu$ also traces the user's
transactions.)

	\end{enumerate}

	{Increment Epoch.} On input from the ledger operator, who advances the
epoch on a schedule fixed in $\crs$ and known to all participants:

	\begin{enumerate}

	  \item Set $\epoch_{\mathsf{cur}} \leftarrow \epoch_{\mathsf{cur}} +
1$.

	\end{enumerate}

      \end{multicols} \end{framed}

    \caption{Ideal functionality for tracing, $\idfu{\algo{trace}}$.}
\label{fig:ftxid} \end{minipage} \end{figure}
